\documentclass[11pt]{article}

\usepackage[a4paper,margin=1in]{geometry}
\usepackage{amsmath,amssymb}
\usepackage{booktabs}
\usepackage{array}
\usepackage{enumitem}
\usepackage{xcolor}
\usepackage{titlesec}
\usepackage{fancyhdr}
\usepackage[colorlinks=true,linkcolor=accent,urlcolor=accent]{hyperref}

% ---- palette -------------------------------------------------------------
\definecolor{accent}{HTML}{1F4E79}
\definecolor{rulegray}{HTML}{888888}
\definecolor{boxbg}{HTML}{F2F5F9}

% ---- section styling -----------------------------------------------------
\titleformat{\section}{\Large\bfseries\color{accent}}{\thesection}{0.6em}{}
\titleformat{\subsection}{\large\bfseries}{\thesubsection}{0.6em}{}
\titlespacing*{\section}{0pt}{1.4em}{0.5em}

% ---- header / footer -----------------------------------------------------
\pagestyle{fancy}
\fancyhf{}
\renewcommand{\headrulewidth}{0.4pt}
\renewcommand{\headrule}{\hbox to\headwidth{\color{rulegray}\leaders\hrule height \headrulewidth\hfill}}
\lhead{\color{rulegray}\small PID Tuner --- Methods Reference}
\rhead{\color{rulegray}\small\thepage}

% ---- a "key formula" box -------------------------------------------------
\newcommand{\keybox}[1]{%
  \par\vspace{0.4em}%
  \noindent\colorbox{boxbg}{%
    \begin{minipage}{\dimexpr\linewidth-2\fboxsep\relax}%
      \vspace{0.3em}#1\vspace{0.3em}%
    \end{minipage}}%
  \par\vspace{0.4em}%
}

\newcommand{\applies}[1]{\smallskip\noindent\textbf{Applies when:} #1\par}
\newcommand{\refuses}[1]{\noindent\textbf{Refuses / fails when:} #1\par}
\newcommand{\inputs}[1]{\noindent\textbf{Inputs:} #1\par}

\setlength{\parindent}{0pt}
\setlength{\parskip}{0.5em}

% =========================================================================
\begin{document}

\begin{center}
  {\huge\bfseries\color{accent} PID Tuner}\\[0.3em]
  {\Large Tuning \& Identification Methods Reference}\\[0.8em]
  {\color{rulegray}\rule{0.5\linewidth}{0.4pt}}\\[0.6em]
  {\normalsize A precise description of every method implemented in the app,}\\
  {\normalsize matched to the formulas in \texttt{tune.py} and \texttt{identify.py}.}\\[0.4em]
\end{center}

\vspace{0.5em}

This document specifies the nine controller-tuning methods and the three
plant-identification primitives that feed them. Each entry states the
governing equations exactly as coded, the assumptions behind them, when the
method applies, and when it refuses or breaks down. Sources are the course
handouts (HO13 AMIGO, HO14 SIMC, problem session PS4) and the primary
literature cited per method.

\tableofcontents

% =========================================================================
\section{Conventions}

The controller is stored internally in \textbf{parallel form}
\keybox{\[
  D(s) \;=\; K_P \;+\; \frac{K_I}{s} \;+\; K_D\, s ,
\]}
while most tuning rules are published in \textbf{textbook (standard) form}
\[
  D(s) \;=\; K_P\!\left(1 + \frac{1}{T_I s} + T_D s\right).
\]
The two are related by
\keybox{\[
  K_I = \frac{K_P}{T_I}, \qquad K_D = K_P\, T_D ,
  \qquad\Longleftrightarrow\qquad
  T_I = \frac{K_P}{K_I}, \qquad T_D = \frac{K_D}{K_P}.
\]}
A rule that yields $(K_P, T_I, T_D)$ is converted to $(K_P, K_I, K_D)$ on
return; both forms are shown in the result panel. Plant polynomials are stored
in descending powers of $s$ (scipy / MATLAB convention), and dead time $L$
enters as a factor $e^{-Ls}$.

The plant model throughout the identification step is
\textbf{first-order plus dead time} (FOPDT):
\[
  G(s) \;=\; \frac{A\, e^{-Ls}}{\tau s + 1},
\]
with static gain $A$ (also written $K$), time constant $\tau$, and dead
time $L$.

% =========================================================================
\section{Identification primitives}

Four of the six tuning methods (ZN-I, ZN-II, AMIGO, SIMC) operate on a
reduced model rather than the raw transfer function. These primitives produce
that model.

\subsection{FOPDT step fit (28.3\,\% / 63.2\,\% method)}

The plant is simulated against an open-loop step of amplitude $\Delta u$, and
$A,\tau,L$ are read off the response. For a true FOPDT system the response is
\[
  y(t) = A\,\Delta u\left[1 - e^{-(t-L)/\tau}\right], \qquad t \ge L,
\]
which reaches $28.3\%$ of its final change at $t = L + \tau/3$ and $63.2\%$ at
$t = L + \tau$. Solving for the two timestamps $t_{28.3}$ and $t_{63.2}$ gives
\keybox{\[
  \tau = 1.5\,(t_{63.2} - t_{28.3}), \qquad
  L = t_{63.2} - \tau, \qquad
  A = \frac{\Delta y_{\infty}}{\Delta u}.
\]}
Using two crossing times (Sundaresan--Krishnaswamy) is more noise-robust than
the classical Ziegler--Nichols tangent construction, whose slope depends on a
numerical derivative at the inflection point.

\applies{Any plant that settles to a finite steady state under a step.}
\refuses{The step produces no discernible change ($\Delta y_\infty \approx 0$);
a placeholder model is returned and downstream rules will report a degenerate
fit.}

\subsection{Relay-feedback test (Aström--Hägglund, 1984)}

A hysteretic on/off relay of amplitude $h$ drives the plant into a sustained
limit cycle. The describing-function approximation of the relay gives the
ultimate gain and period directly from the cycle:
\keybox{\[
  K_u = \frac{4h}{\pi a}, \qquad
  P_u = \text{period of the steady-state limit cycle},
\]}
where $a$ is half the peak-to-peak amplitude of the output oscillation. The
implementation discretizes the plant, applies the relay with optional
hysteresis, and measures $P_u$ from the mean of the last several half-periods
(zero-crossings) and $a$ from the tail extrema.

\applies{The plant sustains a limit cycle under relay feedback (most
processes with $\ge 180^\circ$ of attainable phase lag).}
\refuses{Fewer than four zero-crossings are detected, or the oscillation
amplitude is zero --- the routine raises and suggests the Bode-based method
instead.}

\subsection{Bode-based ultimate gain (analytical)}

When the closed-form model is available (it always is here, since the user
supplied it), $K_u$ and $P_u$ are found by locating the $-180^\circ$ phase
crossover directly rather than by simulating a relay:
\keybox{\[
  \angle G(j\omega_{180}) = -180^\circ
  \;\Longrightarrow\;
  K_u = \frac{1}{\lvert G(j\omega_{180})\rvert}, \qquad
  P_u = \frac{2\pi}{\omega_{180}}.
\]}
The smallest positive crossover is bracketed on a logarithmic frequency grid
and refined by bisection on a continuous (anchored) phase function to avoid
the $(-\pi,\pi]$ wraparound. This is faster and more accurate than the relay
test.

\applies{The plant phase actually crosses $-180^\circ$ (e.g. third-order or
higher, or any order with dead time).}
\refuses{No finite crossover exists (e.g. a pure first- or second-order
lag without dead time) --- then $K_u$ is infinite and ZN-II is not
applicable; the routine raises and points to ZN-I.}

% =========================================================================
\section{Tuning methods}

\subsection{1.\quad Stable pole cancellation}

The course quick-and-dirty method (PS4). The PID controller carries two zeros,
which are placed exactly on two stable plant poles at $s=-p_1$, $s=-p_2$:
\[
  C(s) = K_D\,\frac{(s+p_1)(s+p_2)}{s}.
\]
Expanding the numerator,
$C(s) = K_D s + (p_1+p_2)K_D + p_1 p_2 K_D / s$, so the parallel gains read off
as
\keybox{\[
  K_D \;\text{(free)}, \qquad
  K_P = (p_1 + p_2)\,K_D, \qquad
  K_I = p_1 p_2\, K_D .
\]}
After cancellation the open loop collapses to $K_D\, N(s)/[\,s\,R(s)\,]$, an
integrator times the leftover plant dynamics. $K_D$ is the single free knob
and sets the closed-loop speed. By default the app cancels the two
\emph{slowest} stable poles (smallest $\lvert\mathrm{Re}\rvert$), which yields
the fastest clean response; a complex-conjugate pair is cancelled together so
the gains stay real.

\applies{The plant has at least two stable poles, and those poles are known /
fixed (not drifting with operating point).}
\refuses{Either chosen pole is in the right half-plane or on the imaginary
axis. RHP cancellation is rejected outright --- imperfect cancellation leaves
a hidden unstable mode, the cardinal sin emphasized in lecture
(``never, ever do RHP cancellations'').}

\subsection{2.\quad Ziegler--Nichols Method I (reaction curve)}

Classical 1942 open-loop rules driven by the FOPDT fit. With reaction rate
$R = A/\tau$ and dead time $L$, the PID row of the ZN-I table is
\keybox{\[
  K_P = \frac{1.2}{R\,L} = \frac{1.2\,\tau}{A\,L}, \qquad
  T_I = 2L, \qquad
  T_D = 0.5L .
\]}
The rules were derived for disturbance rejection and characteristically
overshoot ($\sim 40\%$ on the benchmark) when used for setpoint tracking ---
hence the \emph{Halve gains} toggle (\S\ref{sec:halve}).

\inputs{$A,\tau,L$ from the FOPDT step fit.}
\applies{Self-regulating plants with a clear S-shaped step response.}
\refuses{Degenerate FOPDT fit ($A=0$ or $\tau=0$).}

\subsection{3.\quad Ziegler--Nichols Method II (ultimate gain)}

Closed-loop rules driven by $K_u, P_u$ (from the relay test or the Bode
crossover). The PID row is
\keybox{\[
  K_P = 0.6\,K_u, \qquad
  T_I = 0.5\,P_u, \qquad
  T_D = 0.125\,P_u .
\]}
Like ZN-I, this is a disturbance-rejection design and overshoots heavily on
tracking (the most aggressive of the six on the benchmark); the
\emph{Halve gains} toggle tempers it.

\inputs{$K_u, P_u$ from \S2.2 or \S2.3.}
\applies{Plants that admit a finite ultimate gain (phase crosses
$-180^\circ$).}
\refuses{No finite $K_u$ (see \S2.3).}

\subsection{4.\quad AMIGO (Aström--Hägglund, 2006)}

Robust rules from handout HO13, tuned to give roughly $20\%$ overshoot under a
constraint on the peak sensitivity --- more conservative than ZN. For the
non-integrating FOPDT $G(s) = A e^{-Ls}/(\tau s + 1)$:
\keybox{\[
  K_P = \frac{1}{A}\!\left(0.2 + 0.45\,\frac{\tau}{L}\right), \quad
  T_I = \frac{0.4L + 0.8\tau}{L + 0.1\tau}\,L, \quad
  T_D = \frac{0.5\tau}{0.3L + \tau}\,L .
\]}
For an integrating process $G(s) = A e^{-Ls}/s$ the rules switch to
\[
  K_P = \frac{0.45}{A}, \qquad
  T_I = 8L, \qquad
  T_D = 0.5L .
\]

\inputs{$A,\tau,L$ from the FOPDT fit (plus an integrating/self-regulating
flag).}
\applies{FOPDT or integrating-plus-dead-time plants where a robust,
moderate-overshoot response is wanted.}
\refuses{Zero static gain ($A = 0$).}

\subsection{5.\quad SIMC (Skogestad, ``Simple IMC'')}

Internal-model-control rules from handout HO14 with a single tuning knob, the
desired closed-loop time constant $\tau_c$ (default $\tau_c = L$, the
``fast and robust'' choice). For FOPDT this gives a \textbf{PI} controller:
\keybox{\[
  K_P = \frac{1}{A}\cdot\frac{\tau}{\tau_c + L}, \qquad
  T_I = \min\bigl(\tau,\; 4(\tau_c + L)\bigr), \qquad
  T_D = 0 .
\]}
The $\min$ in $T_I$ is Skogestad's anti-sluggishness fix: it caps the integral
time so disturbance rejection on lag-dominant plants does not become
unacceptably slow. Supplying a second time constant $\tau_2$ from a
second-order (SOPDT) fit promotes it to a \textbf{PID} controller:
\[
  K_P = \frac{1}{A}\cdot\frac{\tau_1}{\tau_c + L}, \qquad
  T_I = \min\bigl(\tau_1,\; 4(\tau_c+L)\bigr), \qquad
  T_D = \tau_2 .
\]

\inputs{$A,\tau,L$ from the FOPDT fit; optional $\tau_2$ and $\tau_c$.}
\applies{Lag-dominant processes; the default produces the smoothest,
lowest-overshoot tracking of the six methods on the benchmark.}
\refuses{Zero static gain ($A = 0$).}

\subsection{6.\quad Boyd convex--concave design}

From Hast, Aström, Bernhardsson \& Boyd, \emph{PID Design by Convex--Concave
Optimization} (ECC 2013). Unlike the others, this operates on the
\textbf{true plant frequency response}, so it handles arbitrary-order plants
without an FOPDT proxy. The loop transfer
$L(j\omega) = \bigl(K_P + K_I/(j\omega) + K_D\, j\omega\bigr)\,P(j\omega)$ is
\emph{affine} in $(K_P,K_I,K_D)$ at each frequency. The design maximizes the
integral gain subject to robustness constraints:
\keybox{\[
  \max_{K_P,K_I,K_D}\; K_I
  \quad\text{s.t.}\quad
  \lVert S\rVert_\infty \le M_s, \;\;
  \lVert T\rVert_\infty \le M_t ,
\]}
where the sensitivity peak $M_s$ and complementary-sensitivity peak $M_t$
(both default $1.4$) become forbidden circles in the Nyquist plane:
\[
  \lvert L(j\omega) - c\rvert \ge r, \qquad
  \begin{cases}
    c = -1, & r = 1/M_s \quad(\text{sensitivity})\\[4pt]
    c = -\dfrac{M_t^2}{M_t^2-1}, & r = \dfrac{M_t}{M_t^2-1}\quad(\text{comp. sens.})
  \end{cases}
\]
The constraint $\lvert L - c\rvert \ge r$ is concave; linearizing it about the
current iterate turns each step into a linear program. Iterating to
convergence yields the gains. The solver is seeded from SIMC (a good seed
matters, since the LP bounds are scaled to it).

\inputs{The full plant transfer function, $M_s$, $M_t$, and an optional
derivative on/off flag.}
\applies{Any stable plant; especially valuable for high-order or
resonant plants where an FOPDT reduction throws away too much.}
\refuses{The LP becomes infeasible (constraints too tight for the plant) ---
the iteration stops at the last feasible iterate.}

\subsection{7.\quad Cohen--Coon (1953)}

Open-loop reaction-curve rules, like ZN-I but using the time constant $\tau$
and dead time $L$ separately (ZN-I collapses both into the single slope
$R = A/\tau$). A built-in dead-time correction makes them noticeably better
than ZN-I on delay-dominant plants (large $L/\tau$). With $r = L/\tau$ the PID
row is
\keybox{\[
  K_P = \frac{1}{A}\,\frac{\tau}{L}\!\left(\frac{4}{3} + \frac{r}{4}\right),
  \quad
  T_I = L\,\frac{32 + 6r}{13 + 8r},
  \quad
  T_D = \frac{4L}{11 + 2r}.
\]}
Like ZN, the design targets a quarter-amplitude-decay load response, so it is
aggressive on setpoint steps (largest overshoot of the FOPDT methods on the
benchmark).

\inputs{$A,\tau,L$ from the FOPDT step fit.}
\applies{Self-regulating plants, particularly when dead time is a sizeable
fraction of the time constant.}
\refuses{Degenerate FOPDT fit ($A=0$ or $\tau=0$).}

\subsection{8.\quad Chien--Hrones--Reswick (CHR, 1952)}

A refinement of ZN-I that lets the designer choose two things the older rules
fix implicitly: whether the loop is tuned for \emph{setpoint tracking} (servo)
or \emph{load rejection} (regulator), and whether to target the quickest
\emph{aperiodic} response ($0\%$ overshoot) or the quickest response with
$20\%$ overshoot. The four PID rows, all driven by the FOPDT fit:
\keybox{\[
\renewcommand{\arraystretch}{1.4}
\begin{array}{l|ccc}
\textbf{variant} & K_P & T_I & T_D \\ \hline
\text{setpoint, } 0\%  & 0.60\,\dfrac{\tau}{AL} & \tau       & 0.50\,L \\
\text{setpoint, } 20\% & 0.95\,\dfrac{\tau}{AL} & 1.40\,\tau & 0.47\,L \\
\text{load, } 0\%      & 0.95\,\dfrac{\tau}{AL} & 2.40\,L    & 0.42\,L \\
\text{load, } 20\%     & 1.20\,\dfrac{\tau}{AL} & 2.00\,L    & 0.42\,L
\end{array}
\]}
A caveat worth stating in teaching: the ``$0\%$ / $20\%$ overshoot'' labels
describe the response of the \emph{nominal FOPDT model} from which the rules
were derived. Applied to a true higher-order plant, the realized overshoot can
differ from the label (sometimes the $20\%$ row is actually \emph{less}
oscillatory than the $0\%$ row, because the FOPDT reduction distorts the
apparent dead time). This mismatch is itself a useful demonstration of how
much the model-reduction step matters.

\inputs{$A,\tau,L$ from the FOPDT fit, plus the servo/regulator and
$0\%/20\%$ selections.}
\applies{Self-regulating plants where the designer has a clear preference for
tracking vs. rejection, or for damping vs. speed.}
\refuses{Degenerate FOPDT fit, or an unknown variant request.}

\subsection{9.\quad Tyreus--Luyben (Luyben \& Luyben, 1997)}

The ultimate-gain rule popularized in Luyben \& Luyben, \emph{Essentials of
Process Control} (McGraw-Hill, 1997). Like ZN-II it works from $K_u$ and
$P_u$, but it is deliberately detuned --- a smaller proportional gain and a
much longer integral time --- for larger stability margins and far less
oscillation, at the cost of speed. This suits the long-dead-time and
near-integrating loops common in process control. The PID and PI rows:
\keybox{\[
  \text{PID:}\quad
  K_P = \frac{K_u}{2.2},\quad T_I = 2.2\,P_u,\quad T_D = \frac{P_u}{6.3};
  \qquad
  \text{PI:}\quad
  K_P = \frac{K_u}{3.2},\quad T_I = 2.2\,P_u.
\]}
Because the rule is already conservative, the \emph{Halve gains} toggle is
neither needed nor recommended here.

\inputs{$K_u, P_u$ from the Bode crossover or relay test (\S2.2--2.3); PID/PI
selector.}
\applies{Plants with a finite ultimate gain, especially dead-time-dominant or
integrating loops where ZN-II would be far too oscillatory.}
\refuses{No finite $K_u$ (see \S2.3).}

% =========================================================================
\section{Post-processing: halve gains}
\label{sec:halve}

The ZN rules target disturbance rejection and overshoot on setpoint steps. The
Emami / Aström recommendation is to soften them; the app exposes a toggle that
halves all three gains:
\[
  (K_P, K_I, K_D) \;\longmapsto\; \tfrac12 (K_P, K_I, K_D).
\]
Strictly, the classical recommendation halves $K_P$ alone; halving all three
is a slightly different (and simpler) softening. Since the gains are a free
knob anyway, the override is offered for any method, not just ZN.

% =========================================================================
\section{Benchmark behavior}

All methods are exercised in the test suite against
$G(s) = 1000 / [(s+1)(10s+1)]\cdot e^{-0.5s}$. Representative step-tracking
results:

\begin{center}
\renewcommand{\arraystretch}{1.25}
\begin{tabular}{l r r r}
\toprule
\textbf{Method} & \textbf{Overshoot} & \textbf{$t_s$ (2\%)} & \textbf{IAE}\\
\midrule
Pole cancellation        & $\approx 0\%$  & 34   & 9.95\\
Ziegler--Nichols I       & $41\%$         & 14.7 & 3.68\\
Ziegler--Nichols II      & $78\%$         & 10.7 & 3.51\\
AMIGO                    & $19\%$         & 20.2 & 4.96\\
SIMC                     & $4\%$          & 11.3 & 3.44\\
Boyd ($M_s=M_t=1.4$)     & $18\%$         & 19.2 & 4.89\\
Cohen--Coon              & $44\%$         & 8.2  & 3.22\\
CHR setpoint, $0\%$      & $5\%$          & 20.7 & 3.82\\
CHR setpoint, $20\%$     & $2\%$          & 4.8  & 2.56\\
CHR load, $0\%$          & $35\%$         & 15.2 & 3.84\\
CHR load, $20\%$         & $45\%$         & 13.6 & 3.68\\
Tyreus--Luyben           & $13\%$         & 16.0 & 2.51\\
\bottomrule
\end{tabular}
\end{center}

The pattern matches each method's intent: pole cancellation, SIMC, and the CHR
setpoint rows are smooth; the ZN and Cohen--Coon rows and the CHR \emph{load}
rows are fast but oscillatory (the load rows are tuned for disturbance
rejection, not tracking, so they overshoot a setpoint step); AMIGO and Boyd
land near their $\sim 20\%$ robustness targets; and Tyreus--Luyben sits between
ZN-II and the robust methods, as its conservative design intends.

% =========================================================================
\section{Comparing methods: objective criteria}

No single number declares one method ``best'' --- tuning is multi-objective,
and the ranking depends on the plant \emph{and} on what is being asked of the
loop. The app's \textbf{Compare all methods} view computes the following
metrics for every method and presents them as a colour-coded heatmap (green =
best, red = worst, per column) and a radar chart (each axis normalized so the
outer edge is best). Every metric below is oriented so that \emph{lower is
better}.

\subsection*{Tracking performance (unit setpoint step)}
\begin{itemize}[leftmargin=1.4em,itemsep=0.15em]
  \item \textbf{Overshoot} and \textbf{2\% settling time} --- the classic
        step-response shape numbers.
  \item \textbf{IAE} $=\int_0^\infty |e(t)|\,dt$, the integral of absolute
        error. Two relatives are worth knowing: \textbf{ISE}
        $=\int e^2\,dt$ punishes large early errors and tends to produce fast,
        oscillatory tunings; \textbf{ITAE} $=\int t\,|e|\,dt$ punishes errors
        that \emph{persist}, rewarding well-damped responses and usually
        tracking subjective ``looks good'' best.
\end{itemize}

\subsection*{Disturbance rejection (unit load step at the plant input)}
\begin{itemize}[leftmargin=1.4em,itemsep=0.15em]
  \item \textbf{IAE\textsubscript{load}} $=\int |y(t)|\,dt$ with the setpoint
        held at zero and a step injected at the plant input. Most classical
        rules (ZN, Cohen--Coon, the CHR \emph{load} rows) were designed for
        this objective, so judging them only on setpoint tracking is unfair;
        the two columns routinely disagree about the winner.
\end{itemize}

\subsection*{Robustness (loop frequency response $L = C\,P$)}
\begin{itemize}[leftmargin=1.4em,itemsep=0.15em]
  \item \textbf{Maximum sensitivity} $M_s=\max_\omega |S(j\omega)|
        = \max_\omega \bigl|1/(1+L)\bigr|$. This is the single most-recommended
        robustness scalar: it is the inverse of the shortest distance from the
        Nyquist curve to $-1$, bundling gain and phase margin into one number.
        The usual ``robust enough'' band is $M_s\in[1.4, 2.0]$ ($1.4 \approx$
        conservative, $2.0 \approx$ aggressive).
  \item \textbf{Complementary sensitivity peak} $M_t=\max_\omega|T(j\omega)|$,
        and the classical \textbf{gain} and \textbf{phase margins}, reported
        alongside.
\end{itemize}

\subsection*{Actuator effort / smoothness}
\begin{itemize}[leftmargin=1.4em,itemsep=0.15em]
  \item \textbf{Total variation} $\mathrm{TV}(u)=\sum_k |u_{k}-u_{k-1}|$, a
        measure of how much the control signal moves --- high values mean
        chattering or noise amplification and real wear on a physical
        actuator. The derivative-heavy methods (ZN-II, Cohen--Coon,
        Tyreus--Luyben) tend to score worst here. Peak $|u|$ is also reported.
\end{itemize}

\subsection*{Making ``best'' well-posed}
The principled way to compare is \emph{constrained}: minimise a performance
index subject to a robustness bound --- e.g.\ minimise IAE\textsubscript{load}
subject to $M_s \le 1.6$. This is exactly what AMIGO and Boyd do internally,
and it converts the ill-posed ``which is best?'' into the well-posed ``which is
best at a fixed robustness level?'' The honest answer is then not a single
winner but a \textbf{Pareto frontier}: plotting each method in the
$(M_s,\ \text{IAE}_{\text{load}})$ plane, the non-dominated set traces the
achievable trade-off, and any method up-and-to-the-right of another is strictly
dominated. The radar chart shows all axes at once (bigger polygon = better all
round); the heatmap shows the raw numbers; the Pareto view isolates one
head-to-head trade-off. The plant decides the ranking: dead-time-dominant
plants favour SIMC and Tyreus--Luyben, clean lag-dominant plants favour
pole-cancellation and the CHR setpoint rows, and resonant or high-order plants
favour Boyd.

\addcontentsline{toc}{section}{References}

\begin{enumerate}[leftmargin=1.4em,itemsep=0.2em]
  \item J.\,G. Ziegler and N.\,B. Nichols, ``Optimum settings for automatic
        controllers,'' \emph{Trans. ASME}, 1942.
  \item K.\,J. Aström and T. Hägglund, ``Automatic tuning of simple
        regulators with specifications on phase and amplitude margins,''
        \emph{Automatica}, 1984 (relay method).
  \item K.\,J. Aström and T. Hägglund, \emph{Advanced PID Control}, ISA,
        2006 (AMIGO) --- course handout HO13.
  \item S. Skogestad, ``Simple analytic rules for model reduction and PID
        controller tuning,'' \emph{J. Process Control}, 2003 (SIMC) ---
        course handout HO14.
  \item M. Hast, K.\,J. Aström, B. Bernhardsson and S. Boyd, ``PID design by
        convex--concave optimization,'' \emph{European Control Conference},
        2013.
  \item G.\,H. Cohen and G.\,A. Coon, ``Theoretical consideration of retarded
        control,'' \emph{Trans. ASME} 75, 827--834, 1953.
  \item K.\,L. Chien, J.\,A. Hrones and J.\,B. Reswick, ``On the automatic
        control of generalized passive systems,'' \emph{Trans. ASME} 74,
        175--185, 1952.
  \item W.\,L. Luyben and M.\,L. Luyben, \emph{Essentials of Process Control},
        McGraw-Hill, 1997 (Tyreus--Luyben rules; orig. B.\,D. Tyreus and
        W.\,L. Luyben, \emph{Ind. Eng. Chem. Res.} 31, 2625--2628, 1992).
  \item Course problem session PS4 --- stable pole cancellation.
\end{enumerate}

\end{document}
